\documentclass{beamer}

\usepackage[utf8]{inputenc}
\usepackage[portuguese]{babel}
\usepackage{amsmath, amsfonts, amssymb}

\usetheme{Madrid}

\title{Introdução aos Processos Estocásticos}
\subtitle{Martingais em Tempo Discreto}
\author{Baseado em Privault (2018)}
\date{}

\begin{document}

\frame{\titlepage}

%------------------------------------------
\begin{frame}{Motivação}
  \begin{itemize}
    \item Processos estocásticos descrevem sistemas que evoluem ao longo do tempo sob incerteza.
    \item Exemplos em cursos introdutórios:
    \begin{itemize}
        \item Caminhadas aleatórias.
        \item Filas.
        \item Processos de Markov.
        \item Modelagem de apostas ou tomada de decisão sob incerteza.
    \end{itemize}
    \item Nesta aula: introdução ao conceito de \textbf{martingal}.
    \item Martingais capturam matematicamente a ideia de “jogo justo”.
  \end{itemize}
\end{frame}

%------------------------------------------
\begin{frame}{Filtração: informação ao longo do tempo}
  Em processos estocásticos, a noção de informação é essencial.

  \begin{block}{Definição (Filtração)}
    Uma \textbf{filtração} é uma sequência crescente de $\sigma$-álgebras:
    \[
      \mathcal{F}_0 \subset \mathcal{F}_1 \subset \mathcal{F}_2 \subset \cdots
    \]
    Cada $\mathcal{F}_n$ representa tudo o que é possível saber até o tempo $n$.
  \end{block}

  \begin{itemize}
    \item Se $(X_n)$ é um processo, então:
    \[
      \mathcal{F}_n = \sigma(X_0,\ldots,X_n).
    \]
    \item Em cursos introdutórios, pense em: “histórico do processo até o tempo $n$”.
  \end{itemize}
\end{frame}

%------------------------------------------
\begin{frame}{Exemplo: Caminhada Aleatória}
Considere $S_n = X_1 + \cdots + X_n$, onde $X_k \in \{-1,1\}$ com probabilidade 1/2.

\begin{itemize}
    \item $\mathcal{F}_0$: nenhuma informação.
    \item $\mathcal{F}_1$: sabemos o valor de $X_1$.
    \item $\mathcal{F}_2$: sabemos \emph{ambos} $X_1$ e $X_2$.
\end{itemize}

\begin{block}{Interpretação}
No tempo $n$, a filtração contém todo o histórico já observado da caminhada.
\end{block}
\end{frame}

%------------------------------------------
\begin{frame}{Expectativa Condicional}
  \begin{block}{Intuição}
    A expectativa condicional $\mathbb{E}[F \mid \mathcal{F}_n]$ é a melhor previsão de $F$
    usando apenas as informações disponíveis até o tempo $n$.
  \end{block}

  \begin{itemize}
    \item É um conceito essencial para analisar processos que evoluem no tempo.
    \item Satisfaz a regra da torre:
    \[
       \mathbb{E}\!\left[\,\mathbb{E}[F|\mathcal{F}_{n+1}]\,\middle|\,\mathcal{F}_n\right]
       = \mathbb{E}[F|\mathcal{F}_n].
    \]
  \end{itemize}
\end{frame}

%------------------------------------------
\begin{frame}{Martingais: ideia fundamental}
\begin{block}{Ideia inicial}
Um martingal é um processo cujo valor esperado futuro, dado o presente, é igual ao valor atual.
\end{block}

\begin{block}{Definição}
$(Z_n)$ é um martingal se:
\[
   \mathbb{E}[Z_{n+1} \mid \mathcal{F}_n] = Z_n.
\]
\end{block}

\begin{itemize}
    \item Representa um “jogo justo”: não há tendência de ganho ou perda esperada.
    \item Martingais são essenciais em teoria de apostas e em modelos financeiros.
\end{itemize}
\end{frame}

%------------------------------------------
\begin{frame}{Exemplos Importantes}

\textbf{1. Caminhada aleatória com incrementos centrados}  
Se os incrementos têm média zero:
\[
\mathbb{E}[S_{n+1} \mid \mathcal{F}_n] = S_n.
\]

\textbf{2. Processo das projeções condicionais}  
Se $F$ é uma variável integrável e:
\[
X_n = \mathbb{E}[F|\mathcal{F}_n],
\]
então $(X_n)$ é um martingal.

\begin{block}{Importante}
Esse é o tipo mais fundamental de martingal: previsões sucessivamente refinadas.
\end{block}
\end{frame}

%------------------------------------------
\begin{frame}{Propriedade: Expectativa Constante}
Martingais têm expectativa preservada:

\[
   \mathbb{E}[Z_n] = \mathbb{E}[Z_0].
\]

\begin{itemize}
  \item Não existe “aceleração” esperada no processo.
  \item Demonstração usa duas ideias-chave:
  \begin{itemize}
    \item definição de martingal,
    \item propriedade da torre.
  \end{itemize}
\end{itemize}
\end{frame}

%------------------------------------------
\begin{frame}{Tempos de Parada}
Tempos de parada modelam momentos aleatórios de interrupção.

\begin{block}{Definição}
$\tau$ é um tempo de parada se:
\[
   \{\tau \le n\} \in \mathcal{F}_n.
\]
\end{block}

\begin{itemize}
    \item Exemplo típico: primeira vez que um processo atinge certo estado.
    \item Intuição: no tempo $n$, conseguimos responder “já parou?” usando apenas o passado.
\end{itemize}

\end{frame}

%------------------------------------------
\begin{frame}{Teorema do Tempo de Parada (Doob)}
Se $(M_n)$ é um martingal e $\tau$ é um tempo de parada, então o processo parado
\[
   M_{n \wedge \tau}
\]
também é um martingal.

\begin{itemize}
    \item Essencial para estudar estratégias de jogo.
    \item Diz que “parar no meio” não destrói a propriedade de martingal.
\end{itemize}
\end{frame}

%------------------------------------------
\begin{frame}{Resumo da Aula}
\begin{itemize}
    \item Filtrações: informação disponível ao longo do tempo.
    \item Expectativa condicional: previsão ótima baseada no passado.
    \item Martingais: modelo matemático de jogos justos.
    \item Tempos de parada: momentos aleatórios baseados em informação passada.
    \item Teorema de Doob: martingais permanecem martingais após parada.
\end{itemize}
\end{frame}

\end{document}
